\sectionkicker{Source-backed technical notes}
\subsection*{総括 — 2月が残した技術的な転換点: Technical Notes}
本欄は記事本文で圧縮した一次資料上の情報を、比較・再検証しやすい形で整理したものである。時系列、確認済みの主張、留意点、一次資料URLを掲載する。
\medskip
\subsection*{Theme at a glance}
\addcontentsline{toc}{subsection}{Theme at a glance}
\begin{center}
\footnotesize
\begin{tabularx}{\linewidth}{@{}p{0.20\linewidth}p{0.10\linewidth}p{0.14\linewidth}X@{}}
\toprule
資料 & 位置づけ & 種別 & 時系列 \\
\midrule
OpenAI Codex February stack: app, GPT-5.3-Codex, System Card, and Codex-Spark & 補足資料 & モデル更新 & 2026-02-02 (Agent製品公開); 2026-02-05 (モデル公開); 2026-02-05 (System Card公開); 2026-02-12 (研究Preview) \\
Qwen3.5 / Qwen3-Coder-Next February Model Studio rollout & 補足資料 & モデル更新 & 2026-02-16 (地域別モデル公開); 2026-02-20 (国際提供モデル公開) \\
Voxtral Transcribe 2 / Voxtral Realtime & 補足資料 & オープンウェイト & 2026-02-04 (モデル公開) \\
Gemini 3 Deep Think — February 2026 update & 補足資料 & モデル更新 & 2026-02-12 (モデル更新) \\
vLLM v0.15.1 / v0.16.0 & 補足資料 & Framework & 2026-02-04 (プロジェクト公開); 2026-02-25 (プロジェクト公開) \\
\bottomrule
\end{tabularx}
\end{center}
\normalsize

\begin{technicalnote}{OpenAI Codex February stack: app, GPT-5.3-Codex, System Card, and Codex-Spark}{補足資料}
\begingroup
% reader-facing Technical Notes break policy
\widowpenalty=10000
\clubpenalty=10000
\displaywidowpenalty=10000
\begin{tabularx}{\linewidth}{@{}>{\bfseries}p{0.22\linewidth}X@{}}
組織 & OpenAI \\
種別 & モデル更新 \\
時系列 & 2026-02-02 (Agent製品公開); 2026-02-05 (モデル公開); 2026-02-05 (System Card公開); 2026-02-12 (研究Preview) \\
\end{tabularx}
\smallskip
{\bfseries 一次資料から整理したtechnical points}
\begin{itemize}[leftmargin=1.5em,itemsep=0.35em]
\item \textbf{一次情報で確認できる事実}: OpenAIは2026年2月2日にCodex appを、複数のcoding agentと長時間のparallel workを管理するworkspaceとして発表した。
\item \textbf{一次情報で確認できる事実}: OpenAIは2026年2月5日にGPT-5.3-Codexを発表し、codingに加えて長時間のresearch、tool use、complex execution taskを対象とした。
\item \textbf{一次情報で確認できる事実}: GPT-5.3-Codex System Cardは2026年2月5日に公開され、OpenAIのPreparedness Frameworkに基づくcybersecurity評価をmodelへ適用している。
\item \textbf{一次情報で確認できる事実}: OpenAIは2026年2月12日にGPT-5.3-Codex-Sparkを、realtime coding interaction向けに設計したresearch-preview modelとして発表した。
\item \textbf{Vendor claim}: GPT-5.3-Codexが25\%高速、Codex-SparkがCerebras hardware上で1,000 tokens/sを超えるといった数値はvendor-reported performance claimである。
\end{itemize}
\begin{minipage}{\linewidth}
% reader-facing Technical Notes generic boundary/source tail group
{\bfseries 読む際の境界}
\begin{itemize}[leftmargin=1.5em,itemsep=0.35em]
\item \textbf{分析上の留意点}: このEvidence taskは関連する4件のfirst-party itemをまとめているため、product、base model、System Card、specialized Sparkの境界を明示する必要がある。
\item \textbf{分析上の留意点}: performanceおよびbenchmark比較はvendorによる報告である。
\item \textbf{分析上の留意点}: Codex-Sparkのthroughputに関する記述は、専用Cerebras hardwareとresearch-preview条件に依存する。
\item \textbf{分析上の留意点}: cyber capability classificationは、誇張せずSystem Cardに記載されたthreshold/precautionの境界を保持する必要がある。
\end{itemize}
\begin{samepage}
{\bfseries 一次資料}
\begingroup\sloppy
\Urlmuskip=0mu plus 2mu\relax
\begin{itemize}[leftmargin=1.5em,itemsep=0.25em]
\item {\scriptsize\url{https://openai.com/index/gpt-5-3-codex-system-card}}
\item {\scriptsize\url{https://openai.com/index/introducing-gpt-5-3-codex}}
\item {\scriptsize\url{https://openai.com/index/introducing-gpt-5-3-codex-spark}}
\item {\scriptsize\url{https://openai.com/index/introducing-the-codex-app}}
\end{itemize}
\endgroup
\end{samepage}
\end{minipage}
\endgroup
\end{technicalnote}
\medskip

\begin{technicalnote}{Qwen3.5 / Qwen3-Coder-Next February Model Studio rollout}{補足資料}
\begingroup
% reader-facing Technical Notes break policy
\widowpenalty=10000
\clubpenalty=10000
\displaywidowpenalty=10000
\begin{tabularx}{\linewidth}{@{}>{\bfseries}p{0.22\linewidth}X@{}}
組織 & Alibaba Cloud / Qwen \\
種別 & モデル更新 \\
時系列 & 2026-02-16 (地域別モデル公開); 2026-02-20 (国際提供モデル公開) \\
\end{tabularx}
\smallskip
{\bfseries 一次資料から整理したtechnical points}
\begin{itemize}[leftmargin=1.5em,itemsep=0.35em]
\item \textbf{一次情報で確認できる事実}: Alibaba Cloud Model Studioのlifecycle記録では、qwen3.5-plusやQwen3.5-397B-A17Bを含むQwen3.5系列が2026年2月に利用可能となり、一次カタログにはtext・image・video入力とreasoning/coding/agent能力が記載されている。
\item \textbf{一次情報で確認できる事実}: international側のlifecycle記録では、Qwen3-Coder-Nextが2026年2月20日に、multi-turn tool interactionとrepository-level code understandingを対象とするopen-source code modelとして掲載されている。
\item \textbf{Vendor claim}: benchmark能力や効率の優位性に関する主張は、独立再現されない限りvendor claimとして扱う。
\end{itemize}
\begin{minipage}{\linewidth}
% reader-facing Technical Notes generic boundary/source tail group
{\bfseries 読む際の境界}
\begin{itemize}[leftmargin=1.5em,itemsep=0.35em]
\item \textbf{分析上の留意点}: Model Studioのavailabilityと日付はChina/internationalで異なるため、各lifecycle entryの地域範囲を保持する必要がある。
\item \textbf{分析上の留意点}: lifecycle pageは継続更新され後日のglobal rolloutを含み得るため、2月のchronologyは現在のavailabilityではなく日付付きentryを基準にする。
\end{itemize}
\begin{samepage}
{\bfseries 一次資料}
\begingroup\sloppy
\Urlmuskip=0mu plus 2mu\relax
\begin{itemize}[leftmargin=1.5em,itemsep=0.25em]
\item {\scriptsize\url{https://www.alibabacloud.com/help/en/model-studio/newly-released-models}}
\end{itemize}
\endgroup
\end{samepage}
\end{minipage}
\endgroup
\end{technicalnote}
\medskip

\begin{technicalnote}{Voxtral Transcribe 2 / Voxtral Realtime}{補足資料}
\begingroup
% reader-facing Technical Notes break policy
\widowpenalty=10000
\clubpenalty=10000
\displaywidowpenalty=10000
\begin{tabularx}{\linewidth}{@{}>{\bfseries}p{0.22\linewidth}X@{}}
組織 & Mistral AI \\
種別 & オープンウェイト \\
時系列 & 2026-02-04 (モデル公開) \\
\end{tabularx}
\smallskip
{\bfseries 一次資料から整理したtechnical points}
\begin{itemize}[leftmargin=1.5em,itemsep=0.35em]
\item \textbf{一次情報で確認できる事実}: Mistralは2026年2月4日にVoxtral Transcribe 2 familyを発表し、その中にはrealtime streaming speech-to-text modelが含まれる。
\item \textbf{一次情報で確認できる事実}: realtime modelはApache 2.0のopen weightsとして公開され、13言語に対応する4B-parameter streaming ASR modelと説明されている。
\item \textbf{Vendor claim}: launch materialに記載されたlatency、word-error-rate、price/performance比較はvendor claimである。
\end{itemize}
\begin{minipage}{\linewidth}
% reader-facing Technical Notes generic boundary/source tail group
{\bfseries 読む際の境界}
\begin{itemize}[leftmargin=1.5em,itemsep=0.35em]
\item \textbf{分析上の留意点}: 200ms未満とされるlatencyやaccuracy比較はvendor報告であり、記載されたserving configurationに依存する。
\item \textbf{分析上の留意点}: open-weight/license statusは、周辺のhosted service componentすべてがopen sourceであることを意味しない。
\item \textbf{分析上の留意点}: task sourceはMistralのnews indexであり、item-levelの表現は2月のlaunch articleと整合させる必要がある。
\end{itemize}
\begin{samepage}
{\bfseries 一次資料}
\begingroup\sloppy
\Urlmuskip=0mu plus 2mu\relax
\begin{itemize}[leftmargin=1.5em,itemsep=0.25em]
\item {\scriptsize\url{https://mistral.ai/news/}}
\end{itemize}
\endgroup
\end{samepage}
\end{minipage}
\endgroup
\end{technicalnote}
\medskip

\begin{technicalnote}{Gemini 3 Deep Think — February 2026 update}{補足資料}
\begingroup
% reader-facing Technical Notes break policy
\widowpenalty=10000
\clubpenalty=10000
\displaywidowpenalty=10000
\begin{tabularx}{\linewidth}{@{}>{\bfseries}p{0.22\linewidth}X@{}}
組織 & Google DeepMind \\
種別 & モデル更新 \\
時系列 & 2026-02-12 (モデル更新) \\
\end{tabularx}
\smallskip
{\bfseries 一次資料から整理したtechnical points}
\begin{itemize}[leftmargin=1.5em,itemsep=0.35em]
\item \textbf{一次情報で確認できる事実}: Google DeepMindは2026年2月12日、science、research、engineeringの問題を対象とするspecialized reasoning modeとして、更新版Gemini 3 Deep Thinkを発表した。
\item \textbf{一次情報で確認できる事実}: 2月の発表では、Google AI Ultra subscriber向けのGemini app提供とAPI利用のearly-access pathが説明され、広範な一般API availabilityではなかった。
\item \textbf{Vendor claim}: 発表に含まれる科学領域およびbenchmarkの性能結果は、独立再現されない限りvendor claimとして扱う。
\end{itemize}
\begin{minipage}{\linewidth}
% reader-facing Technical Notes generic boundary/source tail group
{\bfseries 読む際の境界}
\begin{itemize}[leftmargin=1.5em,itemsep=0.35em]
\item \textbf{分析上の留意点}: Deep Thinkはspecialized reasoning modeであり、その結果をGemini 3 family全体へ一般化すべきではない。
\item \textbf{分析上の留意点}: 発表時点のAPI availabilityは限定的なearly accessだった。
\item \textbf{分析上の留意点}: 科学benchmarkおよび比較に関する主張はvendorによる報告である。
\end{itemize}
\begin{samepage}
{\bfseries 一次資料}
\begingroup\sloppy
\Urlmuskip=0mu plus 2mu\relax
\begin{itemize}[leftmargin=1.5em,itemsep=0.25em]
\item {\scriptsize\url{https://deepmind.google/blog/}}
\end{itemize}
\endgroup
\end{samepage}
\end{minipage}
\endgroup
\end{technicalnote}
\medskip

\begin{technicalnote}{vLLM v0.15.1 / v0.16.0}{補足資料}
\begingroup
% reader-facing Technical Notes break policy
\widowpenalty=10000
\clubpenalty=10000
\displaywidowpenalty=10000
\begin{tabularx}{\linewidth}{@{}>{\bfseries}p{0.22\linewidth}X@{}}
組織 & vLLM Project \\
種別 & Framework \\
時系列 & 2026-02-04 (プロジェクト公開); 2026-02-25 (プロジェクト公開) \\
\end{tabularx}
\smallskip
{\bfseries 一次資料から整理したtechnical points}
\begin{itemize}[leftmargin=1.5em,itemsep=0.35em]
\item \textbf{一次情報で確認できる事実}: vLLM v0.16.0は2026年2月25日に公開され、release noteではbranch cutが2026年2月8日だったと記載されている。
\item \textbf{一次情報で確認できる事実}: v0.16.0では、streaming-audio/Voxtral realtime対応を持つWebSocket-based Realtime API、Responses API改善、tool-calling修正、unified parallel drafting、大規模なXPU変更が追加された。
\item \textbf{Project claim}: releaseではasync schedulingとpipeline parallelismによりend-to-end throughputが30.8\%、TPOTが31.8\%改善したと報告されているが、これらはproject-reported measurementである。
\item \textbf{一次情報で確認できる事実}: v0.15.1は2026年2月4日に公開され月初のcontextとなるが、2月のsystems signalとしてはv0.16.0の方が強い。
\end{itemize}
\begin{minipage}{\linewidth}
% reader-facing Technical Notes generic boundary/source tail group
{\bfseries 読む際の境界}
\begin{itemize}[leftmargin=1.5em,itemsep=0.35em]
\item \textbf{分析上の留意点}: v0.16.0のbranch-cut日を考慮すると、2月後半のmodel changeがすべて既に統合済みだったことのEvidenceにはできない。
\item \textbf{分析上の留意点}: performance値はproject release noteのmeasurementであり、独立再現していない。
\item \textbf{分析上の留意点}: このseriesはpatch releaseとmajor releaseをまとめているため、編集上はv0.16.0を中心に扱う。
\end{itemize}
\begin{samepage}
{\bfseries 一次資料}
\begingroup\sloppy
\Urlmuskip=0mu plus 2mu\relax
\begin{itemize}[leftmargin=1.5em,itemsep=0.25em]
\item {\scriptsize\url{https://github.com/vllm-project/vllm/releases/tag/v0.15.1}}
\item {\scriptsize\url{https://github.com/vllm-project/vllm/releases/tag/v0.16.0}}
\end{itemize}
\endgroup
\end{samepage}
\end{minipage}
\endgroup
\end{technicalnote}
\medskip
